\documentclass[a4paper,11pt]{ltjsarticle}
\usepackage{amsmath,amssymb}
\usepackage{enumitem}
\usepackage{geometry}
\geometry{margin=25mm}

\title{党勢指数$C_p(T)$の再定義案(階層間の合成方法の見直し)}
\author{}
\date{2026年9月}

\begin{document}
\maketitle

\section{問題: 現行の合成方法は行政区画の粒度に依存する}

現行の$C_p(T)$は、選挙イベント$e$(国会・都道府県議会・都道府県知事・市区町村議会・
市区町村長のいずれか)ごとに$\sqrt{n_eP_e}$(議席数と代表人口の積の平方根、
ペンローズの平方根則)を計算し、期間$T$に投票日があった全イベントについて
これを単純に合計・正規化していた。

$$C_p(T) = \frac{\displaystyle\sum_{e:\,投票日(e)\in T} \sqrt{n_eP_{e}}\cdot\phi_p(e)}
                  {\displaystyle\sum_{e:\,投票日(e)\in T} \sqrt{n_eP_{e}}}$$

平方根は凹関数なので、同じ総人口を少数の大きな単位(47都道府県)に集約するより
多数の小さな単位(1700超の市区町村)に分割した方が$\sum_e\sqrt{P_e}$の合計は
大きくなる($\sqrt{a}+\sqrt{b} > \sqrt{a+b}$、$a,b>0$)。実際、2018年以降の
105か月分の実績を集計すると、市区町村議会が平均52.9\%を占める一方、
都道府県知事は(ほぼ毎月どこかで選挙があるにもかかわらず)平均3.0\%にしか
ならない。これは市区町村議会が政治的に重要だからではなく、日本の行政区画が
市区町村を約1700に分けているという、\textbf{行政区画の粒度に起因する
アーティファクト}である(2026-09にユーザー指摘)。

ペンローズの平方根則は本来「1つの代表機関の中で、選挙区ごとの人口の違いを
均して有権者1人あたりの投票力を揃える」という理論であり、単位の数が固定された
1つの機関の中での重み付けとして正当化される。これを、単位の数も粒度も
全く異なる複数階層(6種類)をまたいで単純合計するのに転用したことが、
上記のアーティファクトの原因である。

\section{再定義案: 階層内は先に合計してから平方根を1回だけ取る}

$L=\{$衆院選, 参院選, 都道府県知事, 都道府県議会, 市区町村長, 市区町村議会$\}$を
6つの階層とする。階層$\ell\in L$・期間$T$について、$E_\ell(T)$を「階層$\ell$で
投票日が$T$に含まれる選挙イベントの集合」とする。

まず階層ごとに、個々のイベントの$n_eP_e$を\textbf{先に合計}する。

$$N_\ell(T) = \sum_{e\in E_\ell(T)} n_eP_e$$

$N_\ell(T)$は「その階層でその期間に実際に投票があった、議席数×代表人口の総量」
であり、同じ総量を何個の自治体に分割しているかには依存しない($\sqrt{}$を
まだ取っていないため)。この$N_\ell(T)$に対して\textbf{階層あたり1回だけ}
平方根を取ったものを、階層の重み$W_\ell(T)$とする(該当イベントが無い階層は
和から除外する)。

$$W_\ell(T) = \sqrt{N_\ell(T)}$$

\textbf{$W_\ell(T)$は政党別のシェアを一切含まない、階層ごとのスカラーの重みに
すぎない}。政党別の情報は、階層内の加重平均シェア$\phi_p^\ell(T)$として別に
計算する。

$$\phi_p^\ell(T) = \frac{\displaystyle\sum_{e\in E_\ell(T)} n_eP_e\,\phi_p(e)}{N_\ell(T)}$$

これは階層$\ell$の中で、個々のイベントを(平方根を取らない)$n_eP_e$そのもので
加重平均した政党別シェアであり、同じ人口をいくつの自治体に分割しても値が
変わらない(加重平均の性質上、単位を分割・統合しても総和と分母がともに
不変なため)。

最後に、階層の重み$W_\ell(T)$で階層間シェア$\phi_p^\ell(T)$を合成する。

$$C_p(T) = \frac{\displaystyle\sum_{\ell\in L} W_\ell(T)\,\phi_p^\ell(T)}
                  {\displaystyle\sum_{\ell\in L} W_\ell(T)}$$

\section{閉じた形への整理}

$W_\ell(T)\phi_p^\ell(T) = \sqrt{N_\ell(T)}\cdot\dfrac{\sum_e n_eP_e\phi_p(e)}{N_\ell(T)}
= \dfrac{1}{\sqrt{N_\ell(T)}}\sum_{e\in E_\ell(T)} n_eP_e\phi_p(e)$なので、
$C_p(T)$は次の閉じた形にまとめられる。

$$C_p(T) = \frac{\displaystyle\sum_{\ell\in L} \frac{1}{\sqrt{N_\ell(T)}}
                  \sum_{e\in E_\ell(T)} n_eP_e\,\phi_p(e)}
                  {\displaystyle\sum_{\ell\in L} \sqrt{N_\ell(T)}}$$

つまり、個々のイベントの生の寄与$n_eP_e\phi_p(e)$を、そのイベントが属する
階層の正規化定数$\sqrt{N_\ell(T)}$で割ってから全階層について合計し、
分母も同じ$\sqrt{N_\ell(T)}$の合計にする、という形になる。階層内では
素朴な$n_eP_e$加重平均(分割の細かさに依存しない)、階層間だけに
ペンローズ的な平方根の縮小($\sqrt{N_\ell(T)}$)をかける、という2段構成である。

\section{性質の確認}

\begin{itemize}
  \item \textbf{行政区画の粒度に依存しない}: 同じ$N_\ell(T)$(総人口×総議席)を
    自治体1つで実現しても100自治体に分割して実現しても、$\phi_p^\ell(T)$も
    $W_\ell(T)$も変わらない。市区町村議会が「自治体数が多いから」というだけの
    理由で指数を支配することはなくなる。
  \item \textbf{階層をまたぐ縮小はペンローズ則のまま}: 6階層それぞれを
    「1つの代表機関」とみなし、その機関が代表する総量$N_\ell(T)$の平方根を
    重みとする。これは元のペンローズ則の適用範囲(1つの機関の中での重み付け)
    に忠実である——階層をまたいだ単純合計をやめ、階層ごとに1回だけ
    平方根を取る形に改めたことで、理論の転用先が正しい範囲に収まる。
  \item \textbf{その期間に該当階層の選挙が無ければ、その階層は分母・分子から
    単純に除外される}(現行設計と同じ)。ある期間に1階層しか選挙が無ければ、
    $C_p(T)$はその階層の$\phi_p^\ell(T)$にそのまま一致する。
  \item \textbf{新たな恣意的パラメータを導入していない}: 階層間の重みを
    「国政選挙を重視する」のように人手で決め打ちするのではなく、あくまで
    実際に選挙があった人口・議席数から機械的に決まる。
\end{itemize}

\section{注意点(未解決)}

この再定義でも、階層間の相対的な重み$W_\ell(T)$は依然として「その月にどの
階層で選挙が集中したか」に左右される——例えば統一地方選のある月は
都道府県議会・市区町村議会の$N_\ell(T)$が跳ね上がり、その月だけ相対的に
地方議会の比重が大きくなる。これは「実際にその月に何が起きたか」を反映した
結果であり、フロー指標としての設計方針(\S1)には整合するが、月によって
階層間の重みが大きく揺れること自体は現行設計と同様に残る。

また、この再定義を採用する場合、105か月分の履歴の再計算・design\_document.tex
\S2の書き直し・既存note記事の記述修正が必要になる(未実施)。
\end{document}
